\documentclass{article}
\usepackage{amsmath,amssymb}
\usepackage[vocab={music}]{leanify}
\newtheorem{lemma}{Lemma}
\begin{document}

\section{Notation}

We write $f(n)$ for the frequency of the $n$th piano key.
We write $s(k,x)$ for the frequency of $k$ semitones above $x$.
We write $c(x,y)$ for the interval in cents from $x$ to $y$.
We write $P$ for the Pythagorean comma.
We write $S$ for the syntonic comma.
We write $r(k,n)$ for the ratio of $k$ steps of the equal temperament with $n$ steps per octave.
We write $m(x)$ for the MIDI note number of the frequency $x$.

\section{Pitch}

\begin{lemma}\label{pitch}
$f(49)=440$, $f(61)=880$, $m(440)=69$ and $c(440,880)=1200$.
\end{lemma}

\begin{lemma}\label{midi}
Let $n$ be an integer with $1\le n\le88$. Then $m(f(n))=n+20$.
\end{lemma}

\section{Tuning}

\begin{lemma}\label{tuning}
$s(12,440)=880$, $r(12,12)=2$, $S=\frac{81}{80}$ and $P>S$.
\end{lemma}

\end{document}
